\chapter{Redshift Windows, Instrument Choice, and Survey Planning}
\label{chap:observing-planning}

\section{Begin with the observable, not the facility}

An observing plan should start with a physical measurement: ionic column,
line-ratio temperature, gas velocity, density, luminosity, spatial gradient,
or population distribution. The line, required precision, and confounding
features then determine wavelength coverage, resolving power, sensitivity,
aperture, and sample design. Beginning with an available instrument and only
later asking what its spectrum measures invites an exposure-time calculation
for the wrong observable.

Cosmological redshift maps rest and observed coordinates through
\begin{equation}
 \lambda_{\rm obs}=(1+z)\lambda_0,\qquad
 \nu_{\rm obs}=\frac{\nu_0}{1+z}.
 \label{eq:redshift-window}
\end{equation}
An instrument spanning $\lambda_{\min}$ to $\lambda_{\max}$ can observe a
line over
\begin{equation}
 z_{\min}=\max\left[0,\frac{\lambda_{\min}}{\lambda_0}-1\right],
 \qquad
 z_{\max}=\frac{\lambda_{\max}}{\lambda_0}-1,
\end{equation}
before accounting for gaps, throughput, atmospheric absorption, detector
edges, and configuration. A line lying formally inside the range can still
be unusable.

\section{Resolution, sampling, and velocity information}

Resolving power and approximate velocity resolution are
\begin{equation}
 R=\frac{\lambda}{\Delta\lambda},\qquad
 \Delta v\simeq\frac{c}{R}.
\end{equation}
This $\Delta v$ describes an instrumental resolution element, not necessarily
the centroid precision. A high signal-to-noise unresolved line can be
centroided more precisely than its width, while intrinsic width requires
deconvolution of the measured line-spread function (LSF). Real LSFs may have
non-Gaussian wings and depend on wavelength, source extent, slit illumination,
detector position, or observing epoch.

Sampling and resolution are distinct. Two or more pixels per resolution
element may sample an LSF, but correlated extraction and resampling mean the
pixels are not independent measurements. For doublets or kinematic
components, compute separation in velocity and simulate the convolved profile
at the anticipated signal-to-noise ratio. A nominal $R$ alone does not prove
that the scientific parameters are identifiable.

Extended-source spectroscopy adds spatial--spectral coupling. A slit can turn
a velocity gradient or off-center illumination into an apparent line shift.
In a slitless or dispersive X-ray spectrum, source extent broadens features
along the dispersion direction. Integral-field spectroscopy removes some
aperture ambiguity but leaves point-spread-function mixing between spaxels.

\section{Flux, luminosity, surface brightness, and exposure time}

For an integrated line flux $F_{\rm line}$,
\begin{equation}
 L_{\rm line}=4\pi D_{\rm L}^{2}F_{\rm line}.
 \label{eq:line-luminosity}
\end{equation}
Flux density transformations require the appropriate $(1+z)$ factor:
\begin{equation}
 F_{\nu,\rm obs}(\nu_{\rm obs})=
 \frac{L_{\nu,\rm em}[(1+z)\nu_{\rm obs}]}
 {4\pi D_{\rm L}^{2}(1+z)}.
\end{equation}
Equivalent width in wavelength units scales as
$W_{\rm obs}=(1+z)W_0$. Surface brightness from a resolved source suffers
Tolman dimming proportional to $(1+z)^{-4}$, apart from bandpass and
evolutionary effects. A high-redshift galaxy can have detectable integrated
flux while its diffuse halo remains inaccessible per resolution element.

For background-dominated counting, signal-to-noise grows roughly as
$t^{1/2}$, but read noise, source Poisson noise, cosmic rays, confusion,
calibration floors, and overheads modify the scaling. Radio sensitivity also
depends on bandwidth, integration time, number of baselines, weighting, and
weather. Official exposure calculators are preferable to scaling a published
observation with different source size and background.

The correct sensitivity unit follows the science. An unresolved-line limit in
erg s$^{-1}$ cm$^{-2}$, a brightness-temperature limit per beam per channel,
and a column-density limit per velocity interval cannot be substituted for
one another. For a Gaussian line, the peak signal changes inversely with
assumed width at fixed integrated flux; every upper limit must state its
adopted width.

\section{Atmospheric and foreground windows}

Ground-based optical and infrared data encounter telluric absorption,
night-sky emission, thermal background, and variable transmission. Redshift
selection can place a target line on an OH skyline or inside a water band even
when it lies within the nominal spectrograph range. A planning calculation
should propagate the target redshift uncertainty and barycentric velocity,
plot all required lines against sky and transmission spectra, and reserve
continuum on both sides.

Radio-frequency interference and Galactic foreground emission are the
analogous frequency-dependent hazards in radio work. At X-ray energies,
detector background, solar-wind charge exchange, Galactic hot gas, and
photoelectric absorption can dominate faint diffuse sources. Space observing
removes the terrestrial atmosphere but not foregrounds, scattered light,
detector systematics, or scheduling constraints.

\section{A current facility map}

The summary below is a planning orientation, not a substitute for the
call-specific handbook. Instrument modes, offered filters, sensitivities, and
status change. Values were checked against official documentation available
in September 2026 and should be rechecked at proposal submission.

\begin{longtable}{@{}p{0.16\textwidth}p{0.23\textwidth}p{0.25\textwidth}p{0.26\textwidth}@{}}
\caption{Representative facilities and line-spectroscopy roles.
Coverage and resolution are deliberately approximate.}
\label{tab:facility-map}\\
\toprule
Facility or mode & Approximate coverage or resolution & Strong use case &
Planning caution\\
\midrule
\endfirsthead
\toprule
Facility or mode & Approximate coverage or resolution & Strong use case &
Planning caution\\
\midrule
\endhead
HST/COS FUV and NUV &
roughly $900$--$3200$ \AA\ across modes; medium-resolution FUV modes commonly
$R\sim10^4$--$2.4\times10^4$ &
high-throughput UV absorption in the low-redshift CGM and IGM &
coverage has detector gaps and lifetime-position constraints; use the current
LSF and sensitivity tables \citep{COSHandbook2026}\\
\addlinespace
HST/STIS &
UV through optical, with low to high spectral resolution by grating &
high-resolution UV/optical work, spatially resolved slits, and wavelengths
outside optimal COS modes &
usually lower UV throughput than COS; slit and extended-source effects matter\\
\addlinespace
JWST/NIRSpec &
$0.6$--$5.3\,\mu$m; prism $R\sim100$, gratings near $R\sim1000$ or $2700$ &
rest-optical diagnostics at high redshift; fixed slits, microshutters, and IFU &
configuration gaps, source centering, spectral overlap, and wavelength-
dependent LSF \citep{NIRSpecDocs2026}\\
\addlinespace
JWST/MIRI &
mid-infrared imaging and spectroscopy; LRS and medium-resolution IFU modes &
warm molecular lines, ionic fine-structure lines, dust, and obscured nuclei &
thermal/background behavior, fringing, and spatially varying IFU response
\citep{MIRIDocs2026}\\
\addlinespace
VLT/UVES &
approximately $300$--$1100$ nm; typical slit resolving power near $4\times10^4$
and higher in narrow configurations &
precise optical absorption, metal systems, and kinematics &
setting-dependent gaps, slit losses, atmospheric cutoff, and tellurics
\citep{ESOInstrumentDocs2026}\\
\addlinespace
VLT/X-shooter &
simultaneous UV to near-IR, approximately $0.3$--$2.48\,\mu$m; moderate $R$ &
broad line inventories and uncertain-redshift transients or galaxies &
resolution changes by arm and slit; telluric and inter-arm calibration\\
\addlinespace
VLT/MUSE &
optical integral-field spectroscopy, approximately $465$--$930$ nm in common
modes; $R$ of a few thousand &
spatially resolved emission, Ly$\alpha$ halos, and blind line searches &
sky residuals, PSF chromaticity, and surface-brightness selection\\
\addlinespace
ALMA &
receiver bands collectively cover much of $35$--$950$ GHz; offered modes are
cycle specific &
CO and other molecules, [C~I], redshifted fine-structure lines, and dust &
atmospheric windows, configuration, maximum recoverable scale, and excitation
\citep{ALMAHandbook2026}\\
\addlinespace
Karl G. Jansky VLA &
centimeter coverage chiefly across $1$--$50$ GHz plus lower-frequency systems &
H~I, radio recombination lines, continuum, polarization, and molecular masers &
radio-frequency interference, array configuration, bandwidth smearing, and
missing short spacings \citep{VLAOSS2026}\\
\addlinespace
XRISM/Resolve &
non-dispersive few-eV-class X-ray spectroscopy over its operating band &
hot-plasma velocities, temperatures, abundances, and line complexes &
use current in-orbit effective-area and gate-valve status, not prelaunch curves
\citep{XRISMPOG2026}\\
\addlinespace
DESI &
thousands of simultaneous fibers over a wide optical band &
large quasar samples, Ly$\alpha$ forest statistics, galaxy redshifts &
fiber aperture, target selection, spectral covariance, and survey masks
\citep{DESIInstrument2016,DESI2025LyA}\\
\addlinespace
SDSS-V &
multi-epoch optical and near-infrared spectroscopy plus local-volume mapping &
Milky Way, black-hole variability, and spatially resolved nearby-galaxy surveys &
subsurvey selection and data-release-specific calibrations
\citep{Kollmeier2017,SDSSVDocs2026}\\
\bottomrule
\end{longtable}

\section{Worked window logic}

For Ly$\alpha$ at $1215.67$ \AA, an optical spectrograph beginning at
$3600$ \AA\ first accesses the line near $z=1.96$. This arithmetic says
nothing about the atmospheric cutoff, quasar brightness, forest path length,
or whether higher-order Lyman lines are covered. For H$\alpha$ at
$6562.8$ \AA, the NIRSpec upper wavelength near $5.3\,\mu$m corresponds
formally to $z\approx7.1$; in practice sensitivity, configuration, line width,
and neighboring [N~II] determine feasibility.

For a line with rest frequency $115.271$ GHz, CO(1--0) enters a $30$ GHz
receiver near $z=2.84$. A detected feature at one frequency is not by itself a
secure rotational identification: adjacent CO transitions at other redshifts,
[C~I], and fine-structure lines can be confused. Wide bandwidth, a continuum
counterpart, a second line, or an external redshift collapses this ambiguity.

Velocity coverage must also include uncertainty. A systemic redshift error of
$300$ km s$^{-1}$ is comparable to an entire narrow-band tuning margin in some
experiments. For outflows or cluster atmospheres, expected wings and member
velocities may require thousands of km s$^{-1}$ around systemic. Baseline for
continuum fitting belongs inside the requested band.

\section{A reproducible planning workflow}

\begin{enumerate}
 \item State the primary physical parameter and the accuracy needed to change
 the scientific conclusion.
 \item List required and corroborating lines with vacuum rest wavelengths or
 rest frequencies, atomic database version, and expected ratios.
 \item Propagate the full redshift probability distribution into observed
 wavelength, not just the best redshift.
 \item Overlay configurations, detector gaps, skylines, telluric bands,
 foregrounds, and known blends.
 \item Choose $R$ by forward-simulating the narrowest components or doublets
 that must be separated.
 \item Predict line flux, width, spatial size, continuum, and background in
 the exact aperture or beam used by the sensitivity calculation.
 \item Model acquisition, slit or shutter losses, dithers, nods, mosaics,
 calibration, overhead, and maximum recoverable scale.
 \item Define the detection statistic, extraction aperture, upper-limit width,
 and treatment of nondetections before inspecting the result.
 \item Include a validation route: a doublet partner, independent redshift,
 second transition, blank sky, standard star, or repeat epoch.
 \item Record software, calibration reference files, ephemeris corrections,
 data release, exposure-calculator version, and all manual masks.
\end{enumerate}

\section{Survey selection is part of the astrophysics}

A targeted program and a survey answer different questions. Target selection
based on UV brightness can exclude dusty galaxies; requiring background
quasars weights CGM sightlines by the quasar luminosity function and angular
availability; fiber collisions and finite apertures reshape environmental
samples; line-flux selection favors high equivalent width. Population
statistics require a selection function $S(\boldsymbol{x})$ that enters the
likelihood or completeness correction.

For an absorption survey, the usable path length depends on wavelength masks,
signal-to-noise, proximity exclusions, and the column and width of a
hypothetical line. An incidence estimate is therefore based on effective path,
not merely number of spectra. For an emission-line survey, inject-and-recover
tests should span size, width, continuum, and wavelength because completeness
is multidimensional.

Large facilities do not eliminate small systematics. DESI has transformed
Ly$\alpha$ forest statistics through hundreds of thousands of spectra, while
SDSS-V adds multi-epoch and spatial dimensions; their power comes from
documented calibration and selection as much as sample size
\citep{DESI2025LyA,SDSSVDocs2026}. A line atlas should link each diagnostic to
both the physical model and the survey transfer function.

\section{The final pre-proposal audit}

Before submission, verify five boundaries: the line is in a high-throughput
region for the full redshift uncertainty; the resolution answers the stated
kinematic or deblending question; the sensitivity uses the correct line width
and source extent; the spatial configuration recovers the relevant scales;
and at least one alternative identification or systematic has a planned test.
Then repeat the check against the current official call. A proposal built
around a static instrument table can become obsolete, but a plan built around
explicit observables and a reproducible forward model remains portable.
